% ============================================================
% CHAPTER 1
% ============================================================

\chapter{SpherePop as Homotopy Theory}
\label{chap:spherepop-homotopy}

\section{Motivation and Overview}

The SpherePop calculus was introduced in preliminary form in
\cite{Flyxion-SpherePop-v1} as a geometric model of computation based
on two primitive operations: \emph{merge}, which combines geometric
regions, and \emph{collapse}, which abstracts internal geometric detail.
In contrast to symbolic models such as lambda calculus, SpherePop
employs continuous, topologically meaningful configurations, thereby
providing a native interface between geometry, computation,
and the lamphrodynamic fields of \rsvpabbr{}.

The guiding principle of this chapter is that configuration spaces in
SpherePop naturally form $\infty$--groupoids, and that merge operations
correspond to homotopy colimits in a suitably chosen model category.
This identification provides both a rigorous categorical foundation
for SpherePop and a conceptual bridge to the derived geometric
structures developed in Volume~I.

\section{Configuration Spaces as $\infty$--Groupoids}

Let $\Config$ denote the category of geometric configurations
(arising, for example, from finite unions of spherical regions in
$\mathbb{R}^3$ or more general manifolds). Morphisms correspond to
continuous deformations (isotopies) between configurations. Viewed
up to homotopy, $\Config$ assembles into an $\infty$--groupoid whose
objects are configurations and whose $k$--morphisms represent
$k$--fold homotopies.

\begin{definition}
A \emph{SpherePop configuration} is a finite collection of embedded
$n$--dimensional spheres $S^n \hookrightarrow M$, together with
lamphrodynamic labels $(\PhiField,\vField,\SField)$ pulled back
from $M$. Two configurations are considered equivalent if they are
related by a lamphrodynamically admissible isotopy, i.e.\ a smooth
homotopy preserving lamphrodynamic boundary conditions.
\end{definition}

\begin{remark}
The lamphrodynamic fields $(\PhiField,\vField,\SField)$ are not merely
decorations: they encode physical or semantic information associated
to each region, and they influence admissible deformations. Thus
$\Config$ is enriched over the category of lamphrodynamic field
configurations, linking computational and physical structure.
\end{remark}

\section{Homotopy Types and Derived Enhancements}

The naive homotopy type of $\Config$ fails to capture the derived
structure needed for lamphrodynamic evolution. We therefore regard
$\Config$ as an \emph{underived shadow} of a derived configuration
stack $\dConfig$, equipped with a $( -1 )$--shifted symplectic
structure inherited from the lamphrodynamic fields. This allows us
to interpret merge operations as homotopy colimits in $\dConfig$,
thereby unifying geometric computation with the derived geometric
core of \rsvpabbr{}.

\begin{theorem}[informal]
Configuration spaces in SpherePop admit a derived enhancement
$\dConfig$ whose mapping spaces are homotopy invariant and which
carries a $( -1 )$--shifted symplectic structure compatible with
lamphrodynamic evolution.
\end{theorem}

A precise construction of $\dConfig$ requires the machinery of derived
stacks and shifted symplectic geometry developed in Volume~I,
especially \cref{chap:shifted-symplectic,chap:plenum-stack}. We will
return to formal details in \cref{sec:merge-homotopy-colim} below.

\section{Merge as Homotopy Colimit}

The central operation of SpherePop is the merge of two configurations
along a common boundary. Formally, let $X$ and $Y$ be configurations
with a shared subconfiguration $Z$. The intuitive merge is the
glued configuration $X \cup_Z Y$. We show that this construction
corresponds to a homotopy colimit in the derived configuration stack
$\dConfig$.

\begin{theorem}[Theorem 24.2, to be proved in full later]
\label{thm:merge-hocolim}
Let $X \leftarrow Z \rightarrow Y$ be a span of configurations with
lamphrodynamically admissible embeddings. Then the SpherePop merge
$X \cup_Z Y$ is canonically equivalent (up to higher homotopy) to the
homotopy colimit of the span in $\dConfig$.
\end{theorem}

The proof of \cref{thm:merge-hocolim} will occupy a substantial portion
of this chapter. The key idea is that lamphrodynamic admissibility
ensures the compatibility of shifted symplectic forms across glueing
boundaries, allowing the derived homotopy colimit to inherit the
appropriate symplectic structure.

\section{Outline of the Chapter}

We proceed as follows:
\begin{enumerate}
\item We construct the derived configuration stack $\dConfig$.
\item We establish lamphrodynamic admissibility criteria.
\item We define homotopy colimits in $\dConfig$.
\item We prove \cref{thm:merge-hocolim} in full generality.
\item We provide worked examples illustrating merge operations.
\end{enumerate}

\section{Open Questions}

The homotopy--theoretic interpretation of SpherePop raises several
questions:
\begin{itemize}
\item What are the necessary and sufficient conditions for lamphrodynamic
admissibility?
\item Does $\dConfig$ admit additional structure beyond $( -1 )$--shifted
symplectic geometry?
\item Can we classify computational complexity in terms of derived
invariants of $\dConfig$?
\end{itemize}
These questions point to deep connections between lamphrodynamics,
homotopy theory, and computation, many of which remain unexplored.

\newpage

